\documentclass{situationist-preprint}

\title{Walton's Slime Case}
\msid{AES-05-b}
\msfamily{AES-05}
\msdiscipline{Aesthetics}
\mstype{Companion manuscript}
\mskeywords{paradox of fiction; quasi-emotion; Walton; make-believe; alief; startle}

\begin{document}
\maketitle

\begin{abstract}
Walton's case of Charles and the green slime is the standard illustration of the paradox of fiction. AES-05-a uses it to argue that Charles does not fear the slime but quasi-fears it. We separate what the case shows from the theory of make-believe in which Walton embedded it. The case supports the quasi-emotional account by itself; it does not require that it be fictional that Charles is afraid. We consider two variants, one in which Charles comes to believe the slime is real and one in which his response is a mere startle, and show that both confirm the role the account assigns to belief.
\end{abstract}

\section{What the case shows}

Charles watches a horror film in which a slime oozes toward the camera; his pulse quickens, his muscles tense, he grips the chair, and afterwards he says he was terrified. He does not believe the slime exists and does not flee \citep{Walton1978}. AES-05-a takes the case to exhibit the physiological profile of fear (Definition~2.1, Assumption~2.3) without the belief that fear requires (Principle~P1), and concludes that Charles quasi-fears the slime (Theorem~3.4).

Walton's own account has two parts. The first is negative: Charles does not fear the slime. The second is positive: Charles is engaged in a game of make-believe with the film, and his quasi-fear makes it fictional in that game that he fears the slime \citep{Walton1990}. The case supports the first part directly. The second part is a further theory of representation, and the quasi-emotional account of AES-05-a does not depend on it: Definition~2.2 of AES-05-a characterizes quasi-fear by its profile and its intentional object, as represented, without any claim about what is fictional.

\begin{proposition}\label{p:indep}
The slime case supports the claim that Charles does not fear the slime independently of the theory of make-believe.
\end{proposition}

\begin{proof}
The case presents the profile of fear (Assumption~2.3 of AES-05-a) together with the absence of protective action. By Assumption~2.4, fear motivates protective action only through the belief that its object exists and threatens; the absence of such action, in a subject who is not paralyzed or restrained, is evidence that the belief is absent; and by Principle~P1 the absence of the belief excludes fear. No premise concerns fictional truth.
\end{proof}

\section{Two variants}

\emph{The escaped slime.} Suppose that halfway through the film Charles is told, falsely but convincingly, that a real substance has escaped from a laboratory next door and is seeping under the cinema doors. His profile may be unchanged, but now he gets up, warns his neighbours and leaves. The difference between the two episodes is a difference in belief, and it makes exactly the difference to action that Assumption~2.4 predicts. The variant confirms that the belief condition marks the line between fear and quasi-fear.

\emph{The mere startle.} Suppose instead that Charles merely flinches when the slime lunges, as he would at any sudden loud image. One might take this to show that his response is reflex rather than emotion. But the original case is not one of mere startle. Charles's response is sustained, directed at the slime's approach, and modulated by the narrative: it rises as the slime nears the camera and subsides when the scene cuts away. These are marks of a state with an intentional object, as Definition~2.2 requires.

\begin{objection}[Alief]
\citet{Gendler2008} argues that responses like Charles's are produced by alief, an automatic representational state that can conflict with belief. If Charles alieves that the slime is dangerous, perhaps he fears it after all.
\end{objection}

\begin{reply}
Alief may explain the onset of the profile, but it does not supply what Principle~P1 requires, which is belief. If alief sufficed for fear, the person who alieves danger on a glass walkway high above a canyon, while believing herself safe, would fear the drop in the full sense and would be disposed to flee it; she is not, and she continues to walk. The alief theorist and the quasi-emotional account can agree on the mechanism; they disagree only if alief is said to make the state fear, and the absence of protective action tells against that.
\end{reply}

\section{Conclusion}

The slime case shows that a subject can exhibit the full profile of fear toward what he knows to be fictional without the belief, and hence without the action, that fear involves. This is all the quasi-emotional account of AES-05-a needs, and it holds whatever one thinks of the theory of make-believe in which Walton first set the case.

\begin{thebibliography}{99}\small
\bibitem[Gendler(2008)]{Gendler2008} Gendler, T.~S. (2008). Alief and belief. \emph{The Journal of Philosophy}, 105(10), 634--663.
\bibitem[Walton(1978)]{Walton1978} Walton, K.~L. (1978). Fearing fictions. \emph{The Journal of Philosophy}, 75(1), 5--27.
\bibitem[Walton(1990)]{Walton1990} Walton, K.~L. (1990). \emph{Mimesis as Make-Believe}. Harvard University Press.
\end{thebibliography}

\end{document}
